\begin{question}
\section*{Theory}

Make sure that you explain all variables when introducing it for the first time.

\begin{enumerate}

\item State the relevant shallow water equations and name the components.

\item Describe what happens after an initial perturbation during the geostrophic adjustment in terms of energy. Write down the analytical expression for the total kinetic and potential energies. State the ratio of the energies changes.

\item State the equations for the sea surface elevation $\eta$ and the velocity $v$ after stationarity (geostrophy) is reached.

\item Discuss the parameters that are influencing the shape and the magnitude of the sea surface elevation $\eta$ and the velocity $v$ in the steady state.

\end{enumerate}


\section*{Experiments}

Set up the model as described in the laboratory description for laboratory 1 and follow the instruction for the experiments. When presenting results in figures, make sure you present them in a suitable way/quality.


\subsection*{Sea surface height and velocities}

\begin{enumerate}

\item What does steady state mean for the experiment?

\item How can you identify in your experiments that steady state is not yet reached?

\item Describe what happens in the first time steps (12 hours) of simulation time. Feel free to include snap shots. (Hint: Think physically but also about the BC in your model.)

\item How do Rossby radius and domain size interfere? Include plots for the sea surface elevations $\eta$ and velocities $v$ from your experiments to support your answer. Make sure that they also compare your simulation results with the theory. (Hint: Modify the Rossby Radius in various ways and conduct several experiments.)

\item Discuss your choise of parameters with linking them to your answers of question 1.iv and 2.iv.

\end{enumerate}


\subsection*{Energetics}

\begin{enumerate}

\item Plot the total energy of your experiments as a function of time. Think of an informative way how to add the theoretical values (steady state, initially) of the total energy to the plot.

\item Plot the ratio of kinetic and change in available potential energy of the model as a function of time. Add the theoretical value of the steady state of this ratio to the plot in an informative way.

\item Give physical interpretations of the results.

\end{enumerate}
\end{question}